\documentclass[a4paper,11pt]{report}
\usepackage[utf8]{inputenc}
\usepackage[T1]{fontenc}
\usepackage[spanish,es-lcroman]{babel}
\usepackage{graphicx}
\usepackage{longtable}
\usepackage{float}
\usepackage{wrapfig}
\usepackage{amssymb}
\usepackage{hyperref}
\usepackage{amsmath}
\usepackage{amsthm}
\usepackage{wasysym}
\usepackage{dsfont}
\usepackage{enumerate}

% *************************************************************************************************************
% Entorno para ejemplos
% https://tex.stackexchange.com/questions/5223/command-for-argmin-or-argmax
% *************************************************************************************************************
\usepackage{mdframed} % Add easy frames to paragraphs
\usepackage{lipsum} % For dummy text
\usepackage{xcolor}
\usepackage{xparse} % Add support for \NewDocumentEnvironment
\definecolor{graylight}{cmyk}{.30,0,0,.67} % define color using xcolor syntax

\newmdenv[ % Define mdframe settings and store as leftrule
  linecolor=graylight,
  topline=false,
  bottomline=false,
  rightline=false,
  skipabove=\topsep,
  skipbelow=\topsep
]{leftrule}

\NewDocumentEnvironment{example}{O{\textbf{Ejemplo:}}} % Define example environment
{\begin{leftrule}\noindent\textcolor{black}{#1}\par}
{\end{leftrule}}


% *************************************************************************************************************
% Definiciones de teoremas, corolarios, lemas
% *************************************************************************************************************
\newtheorem{theorem}{Teorema}
\newtheorem*{theorem*}{Teorema}
\newtheorem{fact}{Proposición}
\newtheorem*{fact*}{Proposición}
\newtheorem{lemma}{Lema}
\newtheorem*{lemma*}{Lema}
\newtheorem{corollary}{Corolario}
\newtheorem*{corollary*}{Corolario}
\newtheorem{definition}{Definición}
\newtheorem*{definition*}{Definición}


% *************************************************************************************************************
% Comandos
% *************************************************************************************************************

% Inserción de imágenes
\newcommand{\img}[2]{
  \begin{center}
  \includegraphics[width=#2\textwidth]{#1}
  \end{center}
}

% arg min
\DeclareMathOperator*{\argmin}{arg\,min }

% norma || ||
\newcommand{\norm}[1]{||{#1}||}

% Abreviaciones para escribir menos
\newcommand{\mprob}{\underset{S\sim \mathcal{D}^m}{P}}
\newcommand{\prob}{\underset{x\sim \mathcal{D}}{P}}
\newcommand{\zprob}{\underset{z\sim \mathcal{D}}{P}}
\newcommand{\dist}{\mathcal{D}}


% *************************************************************************************************************
% Otras configuraciones
% *************************************************************************************************************
\setlength{\parindent}{0pt}
\setlength{\parskip}{1em}
\everymath{\displaystyle}



\begin{document}
% *************************************************************************************************************
% Frontmatter
% *************************************************************************************************************
\pagenumbering{roman}

\makeatletter
  \def\input@path{{./frontmatter/}}
\makeatother


\begin{titlepage}
	\centering
	\includegraphics[width=0.5\textwidth]{./imgs/ugr.png}\par
	\vspace{1cm}
        \rule{\textwidth}{0.3em}\hfill
        {\huge\bfseries Clasificación no balanceada: \par
        Análisis teórico y algoritmos \par}
        \rule{\textwidth}{0.3em}\hfill
	\vspace{1cm}
        {\scshape Proyecto final de Carrera\par}
        {\scshape Ingeniería informática y Matemáticas\par}
        \vfill
        {\bf Autor \par}
	{Ignacio Cordón Castillo \par}
	{\bf Tutor \par}
	{Salvador García López \par}
	\vfill
      	\includegraphics[width=0.4\textwidth]{./imgs/by-nc-sa.png}\par
	{\large \today\par}
\end{titlepage}
\chapter*{Resumen}
  % Resumen en español
\section*{\sc{Palabras clave}}
% Palabras clave
\chapter*{Abstract}
  Machine learning mathematical foundations and algorithms have been one of the most studied areas in both mathematics and 
computer science in the past few years, due to the constant improvement in computer features. Lots of fundamental questions 
arise: what does it mean to extract knowledge out of data?, when can we learn from the data?, how much data do we need?

One of the most researched topics in machine learning is classification problems: given a sequence of elements belonging to
a domain, which have been labeled, how could we learn from them in a way that if new samples arrive we would be able to 
label them according to the knowledge acquired?. We will specially focus on binary classification problems, those which 
only have two possible labels. We will also study, in a more empirical framework, the particular case in which one 
class has more examples than the other one: imbalance classification.

The aim of this work is both to introduce a mathematical theory for machine learning, which has in the PAC learning 
an excellent formalization, and to study and code some of the most recent algorithms that have been published in scientific
journals on the topic of imbalance classification.

\section*{Mathematical introduction}
We introduce concepts suchs as product ring and semiring of sets, and their relationship with $\sigma$-algebras. We give
the notion of measure, without distinction between $\sigma$-algebras and other sets, whereas in classic literature the latter
ones are called premeasures. We also introduce the outer measures, and we advance towards the construction of a $\sigma$-algebra
based on a premeasure over some space, and the extension of that outer measure to a measure. The result which will give us
such extension is Carathéodory's theorem.

This introduction's purpose is to provide us with basic tools to develop the following theory and to answer the question:
given a set of probabilistic spaces, does it exist a product $\sigma$-algebra space in which the product of the sets has
the product of the probabilities as measure?.

This very first part includes demonstrations for Markov and Hoeffding's inequalities, the former one based on the Hoeffding
lemma. Those will be key inequalities in our subsequent progress.

\section*{Machine learning introduction and PAC framework}
We will provide motivations for the machine learning theory, as well as basic definitions such as domain set, label set, 
true labelling, instance generation, training set, hypothesis, hypothesis' error or learning algorithm. The error of an
hypothesis respect to the true labelling function will be defined as the expectation that the two are different. The error over
a training set, as the mean of the number of instances in which it fails.

We will also describe a relation between hypothesis and the sets that those hypothesis classify as one of the classes. Thus,
we will be able to work both with hypothesis or sets. The essential algorithms we will be using will be empirical risk 
minimizers, that is, all those which try to minimize the error over a given training set.

All those concepts will sum up to the very first notion of PAC learning: we will say that a hypothesis collection is 
learnable if we can allways guarantee an small error from one of them with respect to the true labelling function, under
certain restriction over probabilistic confidence. In particular, we will prove that every finite collection of hypothesis
is learnable, and that there exist infinite learnable collections, such as rectangles hypothesis.

The PAC notions will be weakened as much as possible until we arrive to a more flexible definition of learning: APAC 
learning, which is rather similar to PAC learning, with the main difference that we can define our own error function,
and we allow the instances of both classes to suffer from overlaping (an instance could belong to both classes). As
a result, we will show that the concept of APAC implies the PAC concept under certain mild assumptions.


\section*{Uniform learning and further work}

We will define Glivenko-Cantelli classes (uniform classes) of hypothesis, and show that being Glivenko-Cantelli is an
stronger assumption than being APAC. As a result of the work developed under this chapter, we will prove again that finite
classes are learnable, but under a random error function and with worse asymptotic complexity.

In addition to this, we will show by means of the No Free Lunch Theorem, that if we only see half or less of the instances 
of a domain, that domain is not PAC learnable. Moreover, we will relate this theorem with a rather known fact: the 
bias-complexity tradeoff: the more complex the class of functions we are trying to learn, the lower the error and the
harder the learning process.

Vapnik Chervonenkis classes will also be introduced, along with the concept of shattering. This is a pure combinatorics
definition. The Sauer-Shelah lemma will gives us a relationship between the Vapnik Chervonenkis dimension and the number
of possible hypothesis in terms of their image of a training set of fixed size. 

At this stage, we will have all the pieces to prove the fundamental theorem of PAC learning, that connects statistical
concepts, such as Glivenko-Cantelli classes, combinatorial ones (the Vapnik Chervonenkis dimension) and PAC and APAC
concepts.

\section*{Introduction to the imbalance problem}

In order to study some machine learning algorithms, we will weaken our concept of learning until we reach a framework in which
we can easily explain one of data science most studied problems: imbalanced classification. We will describe the suitable
approaches to solve that problem: undersampling and oversampling with and without filtering of instances and cost-sensitive 
framework. Since most of nowadays algorithms tend to understimate the importance of the minority examples, 
we should provide a set of measurements of the quality of an algorithm based not only on the total accuracy over the training
set, but in the level of success over the minority class.

Our main focus will be on oversampling, the creation of synthetic samples belonging to the minority class. We will provide
a description of a classic method: SMOTE. Although it was first described 15 years ago, it is still a reference in the
treatment of imbalanced datasets, and a lot of other algorithms have arised by means of modification of the former one. It will be the case
of our MWMOTE algorithm.

\section*{Description of the coded algorithms}

We will analyze a series of algorithms that are not currently present in any package of \R programming language. All of
them provide an oversampling scheme, except for NEATER. A brief background theory for those algorithms will be provided.
We will start with MWMOTE, which tries to fix some of the issues of SMOTE algorithm, such as creation of new instances 
using noisy ones, or producing instances between two different minoriry class clusters.

RACOG will also be addresed, as an algorithm that approximates discrete distributions using marginal pairwise ones and 
information theory. Once the distribution has been approximated, it will extract samples following a Monte Carlo scheme
called Gibbs Sampling. wRACOG will be a modification of RACOG that aims to improve classification under a certain wrapper
method provided as parameter.

The next algorithm we will code is RWO. RWO finds its foundation on a statistics theorem that we will prove, and which guarantees
us that under asymptotic and generation scheme assumptions, our newly generated instances will preserve the original mean
and variance of the training data set.

PDFOS will be another of the algorithms selected. It is based on multivariate Gaussian kernel density estimations. We will
provide an introduction to the kernel density estimators pointing out their relationship with histograms and we will give
a measure of the error of the estimation, the Mean Integrated Squared Error. PDFOS will try to adapt the bandwith parameter
for a sum of multivariate Gaussian functions, that is, a multiplicative constant for the unbiased covariance of the minority
training samples. A gradient descendent method will be used to optimize such parameter.

The last algorithm we are going to study is NEATER, a filtering algorithm to clean up noisy instances created as a result of
oversampling. NEATER is based on game theory and Nash profile equilibriums, that guarantee that if an instance could play
two possible roles, it will tend to stabilize as one of them, given a series of payoffs based on their neighbours role.
instances.

\section*{Developed software}
An implementation for those algorithms will be made using the \R programming language combined with \texttt{C++} chunks. 
\R language provides a tradeoff between ease of use for the users and speed of the code, the second one will be used to 
speed up our algorithms when possible. Thus, we will develop a package, called \texttt{imbalance}, to fill the lack of implementation of the 
described algorithms in \R.

It is important to note not only the code, but the methodology used to develop the code: under a GPLv2 or later license,
allocated in a public Github repo and with continuous integration to provide unit testing and periodic updates of
its online web documentation page.

\section*{Experiments on small disjuncts}
Finally, we will conduct a simple experiment on small disjuncts problem, which originate as a direct result of class imbalance
problem. The experiment will somehow show that when we overcome the imbalance of the datasets, small disjuncts tend to
vanish.

We propose a simple measure of the small disjuncts which a dataset has: given a tree classifier, we will consider a small
disjunct all those leaves which label less than a given threshold examples. We will also measure the mean size of the 
leaves in terms of examples classified, sometimes called coverage.

\paragraph{Keywords}
Machine learning, PAC learning, imbalanced classification, \textit{oversampling}, \textit{small disjuncts}

\setcounter{tocdepth}{1}
\tableofcontents

\break 

% *************************************************************************************************************
% Mainmatter
% *************************************************************************************************************
\pagenumbering{arabic}
\setcounter{page}{1}


\makeatletter
  \def\input@path{{./chapters/}}
\makeatother


\chapter*{Introducción}
  El material principal en el que se ha basado la sección de matemáticas ha sido \cite{shwartz_understanding_ml}, 
  que proporciona una detallada formalización de los fundamentos matemáticos del Aprendizaje Automático.
\chapter*{Objetivos}

% *************************************************************************************************************
% Parte de matemáticas
% *************************************************************************************************************
\part{Matemáticas}  
  \chapter{Modelización}
    \subsection{Introducción a la probabilidad}
 \begin{frame}\frametitle{$\sigma$ álgebras, anillos y semianillos}
  \begin{definition}
   $\Sigma \subseteq 2^X$ es $\sigma$-álgebra de conjuntos sobre $X$ si se verifica:
   
   \begin{enumerate}[i]
    \item $X \in \Sigma$
    \item Es cerrada para complementarios: $A\in \Sigma$, entonces $A^c = X\setminus A \in \Sigma$
    \item Es cerrada para uniones numerables: $\{A_n\}_{n\in\mathbb{N}} \subseteq \Sigma$, entonces 
    $\underset{n \ge 1}{\bigcup} A_n \in \Sigma$
   \end{enumerate}
   
   Si $\emptyset \in \Sigma$, es semianillo si:
   \begin{enumerate}[i]
    \item $A,B \in \Sigma$ entonces $A\cap B \in \Sigma$
    \item $A,B \in \Sigma$ entonces existen $A_1 \ldots A_n \in \Sigma$ verificándose $A\setminus B = \sum_{i=1}^n A_i$
   \end{enumerate}

   y es anillo si cumple:
   \begin{enumerate}[i]
    \item $A,B \in R$ entonces $A\cup B \in \Sigma$
    \item $A,B \in R$ entonces $A\setminus B \in \Sigma$
   \end{enumerate}
  \end{definition}

  \begin{fact}
  Toda $\sigma$-álgebra es anillo. Todo anillo es semianillo.
  \end{fact}
 \end{frame}
 
 \begin{frame}\frametitle{Medidas, medidas exteriores, probabilidades}
 \begin{definition}[Medida y medida exterior]
  Sea $\mathcal{A} \subseteq 2^X$. Sea $\mu: \mathcal{A} \rightarrow [0, +\infty]$ con $\mu(\emptyset) = 0$. Entonces:
  
  \begin{itemize}
   \item $\mu$ es medida si es $\sigma$-aditiva, esto es, dados $A_n \in \mathcal{A}$ con 
   $\sum_{i=1}^n A_i \in \mathcal{A} \Rightarrow \mu\left(\sum_{i=n}^{+\infty} A_n \right)= \sum_{i=n}^{+\infty} \mu(A_n)$
  
   \item Si $\mathcal{A}=2^X$, $\mu$ es medida exterior sobre $X$ si:
  
    \begin{enumerate}[i]
     \item Monotonía: $A\subseteq B$ entonces $\mu(A) \le \mu(B)$
     \item $\sigma$-subaditividad: $\mu \left(\bigcup_{n=1}^{+\infty} A_n \right) \le \sum_{n=1}^{+\infty} \mu (A_n)$
    \end{enumerate}
  \end{itemize}
 \end{definition}

 \begin{definition}[Espacio de probabilidad]
  Sea $(X, \Sigma, P)$, con $\Sigma$ $\sigma$-álgebra sobre $X$ y $P:\Sigma \rightarrow [0, +\infty]$ medida. 
  Entonces lo llamamos espacio de probabilidad sii $P(\Sigma)\subseteq [0,1]$.
 \end{definition}
 
 \begin{definition}[Distribución de probabilidad]
  Sea $(X, \Sigma, P)$ espacio de probabilidad. Llamamos a la tupla $\dist = (\Sigma,P)$ distribución sobre $X$. 
  Si $\dist$ es distribución sobre $X$, lo notamos $x\sim \mathcal{D}$
 \end{definition}
 \end{frame}
 
 \begin{frame}\frametitle{Teorema de extensión de Carathéodory}
  \begin{definition}
  $\Sigma(\mu^\ast) = \{A\subseteq X: \mu^\ast(T) = \mu^\ast(T\cap A) + \mu^\ast(T\cap A^c), \forall T\subseteq X\}$
  \end{definition}
  \begin{theorem}
   Sea $\mu^\ast : 2^X \rightarrow [0, +\infty]$ medida exterior sobre $X$. Entonces $\Sigma(\mu^\ast)$
   es una $\sigma$-álgebra sobre $X$ y $\mu^\ast_{|\Sigma(\mu^\ast)}$ es una medida sobre $\Sigma(\mu^\ast)$
  \end{theorem}
 
  \begin{theorem}[Teorema de extensión de Caratheodory]
   Sea $S \subseteq 2^X$ un semianillo, $\mu:S \rightarrow [0,+\infty]$ medida en $S$. Entonces existe
   una medida $\overline{\mu}:\Sigma(S) \rightarrow [0,+\infty]$ verificando $\overline{\mu}_{|S} = \mu$
   \label{th:caratheodory}
  \end{theorem}
 
 \begin{proof}
  \begin{itemize}
   \item LLamar $R=(S)$. Probar $\Sigma(S) = \Sigma(R)$.
  
   \item Si existiese $\overline{\mu}$ debería verificar, $\overline{\mu}(T) \le \sum_{n=1}^{+\infty} \overline{\mu}(A_n)$, 
    con $T \subseteq \bigcup_{n\ge 1} A_n$.
  
   \item Definir $\mu^\ast(T):= \inf\left\{\sum_{n=1}^{+\infty} \mu(A_n): \bigcup_{n\ge 1} A_n \supseteq T,
                   A_n\in R \quad \forall n\in\mathbb{N}\right\}$
   \item Probar que $\mu^\ast$ es medida exterior sobre $X$ y usar el teorema anterior.
  \end{itemize}
 \end{proof}
 \end{frame}  
 
 \begin{frame}\frametitle{Espacio de probabilidad producto}
  \begin{theorem}
   Sean $(X_i, \Sigma_i, P_i)$ con $i=1,\ldots, n$ espacios de probabilidad. Entonces existe un espacio probabilidad
   producto.
   \[(X_1 \times \ldots \times X_n, \Sigma_1 \otimes \ldots \otimes \Sigma_n, P)\]
   donde:
   \begin{enumerate}[i]
   \item $\Sigma_1 \otimes \ldots \otimes \Sigma_n = \Sigma\left(\{A_1 \times \ldots \times A_n: A_i\in \Sigma_i\}\right)$
   \item $P(A_1 \times \ldots \times A_n) = \prod_{i=1}^n P_i(A_i)$, con $A_i\in \Sigma_i$ arbitrarios.
  \end{enumerate}
  \end{theorem}
  
  Se puede probar que el espacio probabilístico producto sólo hay uno, pero no usaremos este hecho, aunque conviene tenerlo
  presente. Cuando tengamos $\dist_1 = \ldots = \dist_n$, llamaremos $\dist^n$ al único espacio de probabilidad producto.
 \end{frame}
 
 \begin{frame}\frametitle{Desigualdades de concentración}
 \begin{lemma}[Desigualdad de Markov]
  Sea $W \ge 0$ una variable aleatoria. Entonces para todo $a > 0$ se verifica:
  \[
    P[W \ge a] \le \frac{\expect(W)}{a}
  \]
  \label{ineq:markov}
 \end{lemma}
  
  \begin{lemma}[Desigualdad de Hoeffding]
   Sean $W_1, \ldots W_m$ variables aleatorias i.i.d. (independiente e idénticamente distribuidas), tales que 
   $a_i \le W_i \le b_i$, y sea $\overline{W} = \frac{1}{m} \sum_{i=1}^m W_i$. Entonces dado $\varepsilon > 0$ arbitrario:
   \[
     P\left[\left| \overline{W} - \expect(\overline{W}) \right| \ge \varepsilon \right] \le 2e^{-2m^2 \frac{\varepsilon^2}{\sum_{i=1}^m (b_i-a_i)^2}}
   \]
   
   En el caso particular $a_i = a$ y $b_i = b$ para todo $i=1, \ldots, m$:
   \[
     P\left[\left| \overline{W} - \expect(\overline{W}) \right| \ge \varepsilon \right] \le 2e^{-2m \left(\frac{\varepsilon}{b-a}\right)^2}
   \] 
   \label{ineq:hoeffding}
  \end{lemma}
 \end{frame}
    En este capítulo proporcionamos una introducción y una motivación al \textit{framework} teórico que usaremos.
Daremos definiciones básicas como dominio, conjunto de etiquetas, error empírico, etc. que
usaremos asiduamente en los capítulos posteriores, así como las definiciones de PAC y APAC cognoscibilidad y
describiremos la relación entre ambas.

\subsection{Ejemplo práctico}\label{sec:first-ex}
Pensemos en un ejemplo: tenemos clientes de un banco que quieren solicitar un préstamo, y a todos se les categoriza
el tamaño del préstamo que quieren solicitar(cómo de grande es su importe) y su nivel de ingresos. Estas variables se miden 
en una escala de $0$ a $1$ donde $1$ es el nivel máximo. Queremos etiquetar a cada cliente como $1\equiv$ conceder préstamo o 
$0\equiv$ no concederlo.

Dado un histórico de $m\in \mathbb{N}$ clientes a los que se les concedieron y devolvieron o no préstamos, tenemos una tupla 
$((x_1, y_1), \ldots (x_m, y_m))$ donde $x_i = (x_i^{(1)}, x_i^{(2)}) \in [0,1]^2, y_i \in \{0,1\}$. Asimismo $x_i^{(1)}$ representa el 
tamaño del préstamo solicitado y $x_i^{(2)}$ representa el nivel ingresos mensuales. Tenemos a los clientes clasificados en 
función de si devolvieron los préstamos o no.

Queremos encontrar una función que ofrezca una predicción sobre cualquier cliente, para minimizar posibles pérdidas del banco, es
decir, buscamos una $f:[0,1]^2 \rightarrow \{0,1\}$ que llamaremos predicción.

Asumimos que los datos de los clientes de que disponemos no van a tener una distribución determinada y
van a ser idéntica e independientemente distribuidos (i.i.d.). El histórico siempre va a crecer, y no sabemos cómo va a hacerlo, 
queriendo aprovechar al máximo la información del mismo.

También asumiremos que tenemos una clase de predicciones $H \subseteq \{0,1\}^{[0,1]^2}$, de entre las que existe
una óptima. Por ejemplo, podríamos limitar la clase de predicciones sobre la que buscamos como subrectángulos de $[0,1]^2$, 
esto es, funciones:

\[h_{a,b,c,d} = \mathds{1}_{[a,b]\times[c,d]}, \qquad [a,b]\times [c,d] \subseteq [0,1]^2\]

\img{./imgs/rect-ex.png}{0.85}

También necesitaremos definir una medida de acierto: ¿cómo de buena es una predicción?.

\section{Definiciones básicas}
\label{sec:defs}

Damos unas notaciones/definiciones básicas que utilizaremos de aquí en adelante, en base a lo descrito en \ref{sec:first-ex}:

\begin{itemize}
\item \textbf{Dominio}: $X$. Llamamos \textbf{instancia} a $x\in X$

\item \textbf{Conjunto de etiquetas}: $Y$. Consideramos $Y = \{0,1\}$, lo que nos restringe de momento al paradigma de clasificación binaria. 
En ocasiones también usaremos $Y = \{-1,1\}$ para las etiquetas. Al $1$ se le llama clase positiva, y al $0$ o $-1$ clase negativa.

\item \textbf{Verdadero etiquetado}: \sloppy Asumimos la existencia de una función ${f: X \rightarrow Y}$ 
que proporciona la verdadera etiqueta de todas las instancias.

\item \textbf{Generación de instancias}: \fussy Se tiene $x\sim \mathcal{D} = (\Sigma, \prob)$. La distribución de probabilidad nos da 
información sobre la probabilidad de extraer cada posible instancia desde  $x \in X$. 

\item \textbf{Conjunto de entrenamiento}: $S = ((x_1,y_1), \ldots, (x_m,y_m))$, $S\in (X \times Y)^m$.
Nótese que llamarlo conjunto puede dar lugar a confusión, puesto que se trata de una tupla. También notaremos 
$S_x = (x_1, \ldots x_m)$.

De momento asumiremos que las etiquetas del conjunto de entrenamiento se corresponden con el verdadero etiquetado: 
$y_i = f(x_i)$, por lo que no podemos tener una instancia con etiquetas diferentes.

La elección de $S_x$ es idéntica e independientemente distribuida, esto es $x_i \sim \mathcal{D}$ para todo $i=1, \ldots, m$.
Lo notamos $S_x \sim \mathcal{D}^m$, o $S \sim \mathcal{D}^m$, por abuso de notación. Por lo visto en el capítulo
anterior, en \ref{sec:prob-space}, dicha probabilidad existe.

\item \textbf{Hipótesis/clasificador/predicción}: cada posible aplicación perteneciente a 
$\{h, h:X \rightarrow Y\} := 2^{X}$. 

\item \textbf{Error del clasificador}: Definimos el error del clasificador, suponiendo 
$\{x\in X : h(x) \neq f(x)\} := [h\neq f] \in \Sigma$ como:
\[
  L_{\mathcal{D},f}(h) :=  \prob [h \neq f]
\]

\item \textbf{Algoritmo de aprendizaje}: Llamamos algoritmo de aprendizaje a cualquier aplicación que tome conjuntos de 
entrenamiento y devuelva hipótesis sobre el problema:
\[
  A: \underset{m\in \mathbb{N}}{\bigcup} (X\times Y)^m \rightarrow 2^{X}
\]

Asumimos que el algoritmo no tiene acceso a la función de verdadero etiquetado $f: X \rightarrow Y$ ni a
la distribución $\mathcal{D}$.
\end{itemize}

Es decir, nos centraremos en el problema de la clasificación, que se trata de un aprendizaje supervisado (conocemos las etiquetas para
el conjunto de entrenamiento), \emph{batch} (recibimos todo el conjunto de entrenamiento, y no sucesivas porciones del mismo en
contraposición al aprendizaje por lotes) y pasivo (el algoritmo no tiene interacción con el usuario).

\subsection{Relación entre hipótesis y conjuntos}
\label{biyeccion-canonica} 
En clasificación binaria, existe una biyección canónica entre la clase de hipótesis y el conjunto potencia de $X$, donde
a cada hipótesis se le asigna su clase positiva:
\begin{equation*}
 \begin{array}{rcl} 
  \{h, h:X \rightarrow Y\} & \longrightarrow & \mathcal{P}(X) \\
  h & \longmapsto & X_h := \{x\in X: h(x) = 1\}
 \end{array}
\end{equation*}

Esta relación es biyección claramente, lo que justifica que identifiquemos $\{h, h:X \rightarrow Y\}$ con $2^X$, donde $2^X$ 
suele ser una notación empleada para referirse a $\mathcal{P}(X)$. Esto nos permite trabajar con una aplicación
$h:X \rightarrow Y$ o con su conjunto asociado.

\subsection{Minimización del riesgo empírico}

\begin{definition}[Riesgo empírico]
Definimos el riesgo empírico o error empírico de un clasificador $h:X \rightarrow Y$ como:
\[
  L_S(h) = \frac{|\{i\in {1,\ldots, m}: h(x_i) \neq y_i\}|}{m}
\]
\end{definition}

Es decir, es el error del clasificador $h$ sobre el conjunto de entrenamiento $S$. Nótese que si $[h\neq f]$ como hemos
supuesto al principio, fijado una tamaño de muestra $m$ y $h\in 2^X$, $S \mapsto L_S(h)$ es una variable aleatoria, y por
tanto $L_S(h)^{-1}(]-\infty, a])$ sería un conjunto $\Sigma^m$ medible

\begin{fact}
Sea $H\subseteq 2^X$. Sea $\mathcal{D}$ distribución sobre $\mathcal{X}$. Sea $f \in H$ fija. Entonces para
cualquier $h\in H$:
\[
  \mexpect(L_S(h)) = L_{\dist,f}(h)
\]
\label{fact:expect-error}
\end{fact}

  \begin{proof}
   Llamamos $p= \prob [f\neq h]$.

   Nótese que $h$ puede cometer como mucho $m$ errores.
   \begin{align*}
   \mexpect(L_S(h)) &= \sum_{k=0}^m \frac{k}{m} \binom{m}{k} p^k(1-p)^{m-k}\\
                    &= \sum_{k=1}^m \binom{m-1}{k-1} p^k(1-p)^{m-k} \\
                    &= p\cdot \sum_{k=0}^{m-1} \binom{m-1}{k} p^{k}(1-p)^{m-1-k}\\
                    &= p(1+(1-p))^{m-1} = p
   \end{align*}
  \end{proof}


\begin{definition}[Minimizador del riesgo empírico, ERM]
Decimos que un algoritmo $A: \underset{m\in \mathbb{N}}{\bigcup} (X\times Y)^m \rightarrow 2^{X}$ es un $ERM$ 
(\textit{Empirical Risk Minimizer}) si busca una hipótesis cuyo error empírico sea mínimo:
\[
  A(S) = \argmin_{h\in 2^X} L_S(h)
\]
\end{definition}

Definimos $ERM(S)$ como todas las hipótesis producidas por algoritmos que son ERM
sobre el conjunto $S$. Cuando notemos $\mathcal{Q}(ERM(S))$ con $\mathcal{Q}$ un predicado sobre elementos 
de $2^X$ nos refereriremos a que la propiedad se verifica para todos los elementos de $ERM(S)$.

Trivialmente $L_{\mathcal{D},f}(f) = 0$ y $L_S(f) = 0$ para cualquier $S \in (X \times Y)^m$ conjunto de entrenamiento.
Esto último implica que $A(S) = 0$ para $A$ un algoritmo ERM.

Aunque $A(S)$ minimice el error sobre el conjunto de entrenamiento, esto no significa que el error $L_{\mathcal{D},f} (A(S))$ 
sea mínimo. Lo ilustramos en el siguiente ejemplo:

\begin{example}
Sea $X = \mathbb{R}$, $\mathcal{D}$ la distribución uniforme sobre $[0,2]\subset \mathbb{R}$, y la función
$f= \mathds{1}_{[0,1]}$.

Sea $S = ((x_1,y_1), \ldots (x_m, y_m))$ un conjunto de entrenamiento de tamaño $m$ sin elementos repetidos 
y el clasificador:
\[
  h_S(x) = \left\{\begin{array}{ll}
                   y_i & \textrm{si } \exists i\in \{1\ldots m\} : x=x_i\\
                   0   & \textrm{si } \nexists i\in \{1\ldots m\} : x=x_i
                   \end{array}\right.
\]

No es difícil ver tomando $A(S) = h_S$, $A$ es un ERM. Pero se verifica $\prob[h_S \neq f] = 1/2$. $A(S)$ tiene el 
mismo nivel de acierto que el clasificador idénticamente $1$. A este fenómeno lo denominamos \textit{overfitting}.
\end{example}

\begin{definition}[Overfitting]
 Dado $h: X\rightarrow Y$ clasificador, diremos que produce overfitting sobre el conjunto de entrenamiento 
 $S$ si se cumple $L_S(h) << L_{\mathcal{D},f}(h)$.
\end{definition}

Necesitamos incorporar, por tanto, \textit{conocimiento previo} al problema, de manera que revisitamos la definición de riesgo
empírico.

\subsection{ERM con sesgo inductivo}
Se intenta corregir el riesgo de producir \emph{overfitting} restringiendo el espacio de búsqueda, esto es, la clase de 
hipótesis $H \subseteq 2^X$ desde la que el algoritmo puede escoger un $h: X\rightarrow Y$. Llamamos a esto 
\emph{sesgo inductivo}, puesto que se asume una determinada clase de funciones en función de las 
características del problema y del conocimiento previo que tenemos del mismo.

\begin{definition}[ERM con sesgo inductivo]
Decimos que un algoritmo $A: \underset{m\in \mathbb{N}}{\bigcup} (X\times Y)^m \rightarrow 2^{X}$ es un $ERM$ con sesgo 
inductivo hacia $H$ y lo abreviamos $ERM_H$, si busca una hipótesis en $H$ cuyo error empírico 
sea mínimo, es decir:
\[
  A(S) = \argmin_{h\in H} L_S(h)
\]
\end{definition}

Se define $ERM_H$ y $ERM_H(S)$ de manera análoga a $ERM$ y $ERM(S)$.

Dada $H\in 2^X$, podemos tomar la clase de funciones iguales probabilístamente a $f$:
$\Hclass = \{f\in 2^X: \exists h\in H \textrm{ verificando }\prob[f\neq h] = 0\}$.

\begin{definition}[Hipótesis de factibilidad]

Diremos que $H$ verifica hipótesis de factibilidad respecto a $f\in 2^X$ y $\mathcal{D}$ distribución sobre 
$X$ si $f\in \Hclass$. 
\end{definition}

Por definición: $\prob[f\neq h]=0$ sii $L_{D,f}(h) = 0$ sii $L_{D,h}(f) = 0$. También es inmediato afirmar que
$H \subseteq \Hclass$ para cualquier $\dist$.

\begin{fact}
Dada $H$ verificando factibilidad respecto a $f\in 2^X$ y $\mathcal{D}$ distribución sobre $X$, entonces siendo 
${\bar{h}} \in H$ tal que $L_{D,f}(\bar{h}) = 0$, se verifica, con $A$ algoritmo arbitrario:
\begin{enumerate}[i.]
\item $\mprob \bigg[L_S(\bar{h}) = 0 \bigg] = 1$.
\item $\mprob \bigg[L_S(ERM_H(S)) = 0 \bigg] = 1$.
\item $L_{\dist,f}(A(S)) = L_{\dist,\bar{h}}(A(S))$ para un algoritmo $A$ arbitrario.
\end{enumerate}
\label{fact:factibilidad}
\end{fact}

\begin{proof}
Fijamos $A$ un $ERM_H$. Sabemos que por hipótesis de factibilidad existe $\bar{h} \in H$ tal que
$L_{\mathcal{D},f}(\bar{h}) = \prob[\bar{h}\neq f] = 0$

Entonces: 
\[
  \mprob[L_S(\bar{h}) > 0] = \mprob
  [\exists i: \bar{h}(x_i) \neq f(x_i)] \le \sum_{i=1}^m \prob[h\neq f] = 0
\]

Y como $[L_S(\bar{h}) = 0] \subseteq [L_S(A(S)) = 0]$ por ser $A(S) = \argmin_{h\in \mathcal{H}} L_S(h)$, deducimos:
\[
  1 = \mprob[L_S(\bar{h}) = 0] \le \mprob[L_S(A(S)) = 0] \le 1
\]

Fijamos ahora $A$ un algoritmo. Veamos $L_{\dist,f}(A(S)) = L_{\dist,\bar{h}}(A(S))$:
  \begin{align*}
  \prob[A(S) \neq f] &= \prob\bigg([A(S) \neq f \wedge f=h] \cup[A(S) \neq f \wedge f\neq h]\bigg)\\
  &= \prob[A(S) \neq f \wedge f=h] + \prob[A(S) \neq f \wedge f\neq h]\\
  &= \prob[A(S) \neq h \wedge f=h] + 0 = \prob[A(S)\neq h]
  \end{align*}
\end{proof}

Nótese que para definir $ERM$ y $ERM_H$ necesitamos poder tomar $\argmin$ siempre. Parece una condición fuerte,
pero en nuestro caso, como $L_{S}(f)$ siendo $S$ de tamaño $m$ puede tomar finitos valores $0, \frac{1}{m}, \ldots, \frac{m-1}{m},1$,
el mínimo se alcanza. El valor $L_{\mathcal{D},f}(ERM_H(S))$ depende por tanto del conjunto de entrenamiento $S$, y la elección del
mismo está sometida al azar. Ahí entrará nuestra definición de aprendizaje PAC, midiendo conjuntos probabilísticamente.
    \section{PAC cognoscibilidad}
Proporcionamos a continuación una definición que materializa el concepto de que un algoritmo pueda aprender de los datos. Las
siglas PAC derivan de Probablemente Aproximadamente Correcto. Conviene tener en cuenta para la siguiente definición la
biyección canónica \ref{biyeccion-canonica}, que nos permite asegurar la medibilidad de las hipótesis incluyéndo sus conjuntos
asociados en la $\sigma$-álgebra.

\begin{definition}[PAC cognoscible]
Una clase de funciones $H \subseteq 2^X$ es PAC cognoscible sii existen $A: \underset{m\in \mathbb{N}}{\bigcup} (X\times Y)^m \rightarrow 2^X$ 
y $m_{H} : ]0,1[^2\rightarrow \mathbb{N}$ verificando que para todo
$0 < \varepsilon, \delta < 1$, para toda distribución $\mathcal{D}=(\Sigma, \prob)$ sobre $X$ tal que $H\subseteq \Sigma$
y para toda función de verdadero etiquetado $f\in \Hclass$, dado $m \ge m_H(\varepsilon, \delta)$ entonces:
\[
  \mprob \bigg[L_{\mathcal{D},f}(A(S)) \le \varepsilon \bigg] \ge 1-\delta
\]
\label{def:pac-original}
\end{definition}

$A$ recibe el nombre de \textit{algoritmo de aprendizaje}, $m_H$ se llama \textit{complejidad muestral}, $(1-\delta)$ es la 
\textit{confianza} de la predicción y $(1-\varepsilon)$ la \textit{exactitud}. Estos dos últimos parámetros explican los 
calificativos aproximadamente ($\equiv$ confianza) y correcto ($\equiv$ exactitud).

En lo que sigue supondremos que las distribuciones de las que se habla, $\dist$, incluyen en la $\sigma$-álgebra las propias
hipótesis.

Nótese que la definición de PAC cognoscibilidad implica hipótesis de factibilidad de $H$ respecto a la función de verdadero
etiquetado $f$ y la distribución $\dist$. El último apartado de la proposición \ref{fact:factibilidad} 
nos permite dar una caracterización con $f\in H$ estrictamente:

\begin{charact}[PAC cognoscible]
Una clase de funciones $H \subseteq 2^X$ es PAC cognoscible sii existen una función 
$m_{H} : ]0,1[^2\rightarrow \mathbb{N}$, llamada complejidad muestral, y 
$A: \underset{m\in \mathbb{N}}{\bigcup} (X\times Y)^m \rightarrow 2^X$ verificando que para todo
$0 < \varepsilon, \delta < 1$, para toda distribución $\mathcal{D}$ sobre $X$ y para toda función de 
verdadero etiquetado $f\in H$, dado $m \ge m_H(\varepsilon, \delta)$ entonces:
\[
  \mprob \bigg[L_{\mathcal{D},f}(A(S)) \le \varepsilon \bigg] \ge 1-\delta
\]
\label{def:pac}
\end{charact}

Consideraremos $m_{H}$ única en el sentido de que para cada $(\varepsilon, \delta)$ nos devuelva el menor natural
verificando las hipótesis del enunciado.

\begin{lemma}[Desigualdad fundamental]
 Dados $m\in \mathbb{N}, \varepsilon\in ]0,1[$ se verifica $(1-\varepsilon)^m \le e^{-\varepsilon m}$
 \label{lemma:fund-ineq}
\end{lemma}
\begin{proof}
 Basta hacer la demostración para $m=1$.
 
 Tomando $h(\varepsilon) = e^{-\varepsilon} + \varepsilon - 1$, se tiene que $h$ es derivable y 
 $h'(\varepsilon) = -e^{-\varepsilon} + 1 \ge 0$. Como $h(0) = 0$, se tiene $h \ge 0$.
\end{proof}

\begin{theorem}[Las clases finitas son PAC cognoscibles]
Sea $H \subseteq 2^{X}$ finito. Entonces $H$ es PAC cognoscible con:
\[
  m_H(\varepsilon, \delta) \le \left\lceil \frac{1}{\varepsilon} \log\left(\frac{|H|}{\delta} \right) \right\rceil
\]
\label{th:finitas-pac}
\end{theorem}
  \begin{proof}
   Fijamos $\varepsilon, \delta \in ]0,1[$. Sea una distribución $\mathcal{D}$ sobre $X$, $m\in \mathbb{N}$ y una función de verdadero 
   etiquetado $f\in H$, tomamos la clase de hipótesis ``malas'', que será numerable al ser $|H| < \infty$:
   \[
     H_B = \{h\in H: L_{\mathcal{D},f}(h) > \varepsilon\}
   \]

   Sea $A$ un $ERM_{H}$, entonces:
   \[
     \mprob [L_{\mathcal{D},f}(A(S)) > \varepsilon] \le \mprob 
     [\exists h\in H_B : L_S(h) = 0] \le \sum_{h\in H_B} P [L_S(h) = 0] 
   \]

   La primera desigualdad viene dada por la proposición \ref{fact:factibilidad}. La segunda, por subaditividad, puesto que:
   \[
     \mprob [\exists h\in H_B : L_S(h) = 0] = \mprob \left(\underset{h\in H_B}{\bigcup} [L_S(h) = 0]\right)
   \]

   Además, fijada $h\in H_B$, como $L_{\mathcal{D},f}(h) > \varepsilon$:
   \begin{align*}
   \mprob[L_S(h) = 0] = \mprob[h(x_i) = f(x_i), i =1,\ldots m\}]\\
   = \prod_{i=1}^m \prob[h = f] = \prod_{i=1}^m (1 - L_{\mathcal{D},f}(h)) \le (1-\varepsilon)^m 
   \underset{\textrm{\ref{lemma:fund-ineq}}}{\le} e^{-\varepsilon m}
   \end{align*}

   Las dos desigualdades probadas, junto a la hipótesis del enunciado, y usando $H_B \subseteq H$ dan lugar a:
   \[
     \mprob[L_{\mathcal{D},f}(h_S) > \varepsilon] \le |H|e^{-\varepsilon m}
   \]
   
   Y basta tomar $m\in \mathbb{N}$ tal que $|H|e^{-\varepsilon m} \le \delta$, sii 
   $m\ge \frac{1}{\varepsilon} \log\left(\frac{|H|}{\delta}\right)$.
  \end{proof}

¿Hay ejemplos de clases infinitas PAC cognoscibles? La respuesta es afirmativa. Veamos un ejemplo.

\begin{example}
  \begin{definition}[Clasificadores de rectángulo]
  La clase de clasificadores de rectángulo en el plano se define por:
  \[
    H^2_{rec} = \{ h_{a,b,c,d}: a\le b, c\le d\}
  \]
  donde $h_{a,b,c,d} = \mathds{1}_{[a,b]\times [c,d]}$
  \label{def:rec2}
  \end{definition}
    \begin{fact}
    $H_{rec}^2$ es PAC cognoscible
    \end{fact}

    \begin{proof}
    Haremos la prueba tratando los clasificadores como conjuntos, en virtud de la biyección canónica \ref{biyeccion-canonica}.
    
    Sea $A$ el algoritmo que devuelve el rectángulo más pequeño que engloba a todos los ejemplos positivos de un conjunto 
    de entrenamiento $S$. Claramente $A$ es un $ERM_H$.

    Sea $R= [a,b]\times [c,d]$ el rectángulo que proporciona el verdadero etiquetado. Fijamos $\varepsilon, \delta \in ]0,1[$.
    
    Supongamos $\mathcal{B} = [L_{\dist,R}(A(S)) > \epsilon] \neq \emptyset$. Sea $S\in \mathcal{B}$ y 
    $A(S) = [a^{\ast}, b^{\ast}] \times [c^{\ast}, d^{\ast}]$. Tomamos los rectángulos:
    \begin{align*} 
    R_1^S = [a,b^{\ast}] \times [c,d], \qquad R_2^S = [a^{\ast},b] \times [c,d] \\ 
    R_3^S =[a,b] \times [c,d^{\ast}],   \qquad R_4^S = [a,b] \times [c^{\ast},d]     
    \end{align*}
    verificándose $R_j^S \subseteq R$ para todo $j$. Además debe existir $L_{\mathcal{D},R}(R_j^S) \ge \frac{\varepsilon}{4}$, 
    para algún $j\in \{1, \ldots 4\}$, ya que en caso contrario, por subaditividad tendríamos $L_{\dist,R}(A(S)) < \epsilon$. 
    Además podemos suponer s.p.g. $(R_i^S)^c = R\setminus R_i^S$ puesto que un clasificador $A(S)$ sólo fallará en el
    espacio restante entre $R$ y $A(S)$.
    
    Seleccionamos $R_i \in \{R_i^S: S\in \mathcal{B}\}$, $i=1, \ldots, 4$. suponiendo s.p.g. que existen los cuatro (caso opuesto 
    bastaría hacer la demostración para los rectángulos que sí existiesen). Fácilmente se comprueba que 
    $S\in \mathcal{B} \Rightarrow A(S) \cap R_i^c = \emptyset$ para todo $i=1, \ldots, 4$, usando de nuevo reducción al 
    absurdo y subaditividad.
    
    La demostración acaba probando que:
    \begin{align*}
    \mprob [\exists i : A(S)\cap R_i^c = \emptyset] &\le \sum_{i=1}^4 \mprob[A(S)\cap R_i^c = \emptyset] \\
                                                    &\le 4 \left(1-\varepsilon \right)^m \le 4e^{-\varepsilon m/4}
    \end{align*}

    Juntando toda esta información: 
    \[
      \mprob[L_{\dist,R}(A(S)) > \epsilon] \le \mprob [\exists i : A(S)\cap R_i^c = \emptyset] \le 4e^{-\varepsilon m/4}
    \]
                                    
    y tomando $m > \frac{4}{\varepsilon} log \left( \frac{4}{\delta} \right)$ llegaríamos al resultado buscado.
    \end{proof}
\end{example}

\section{APAC cognoscibilidad}

Nótese que las condiciones exigidas en PAC cognoscibilidad son muy fuertes: cumplir la hipótesis de factibilidad y que 
la hipótesis devuelta deba estar en $H$.Relajaremos esta definición con el concepto de PAC agnóstico (APAC).
Vamos a relajar tanto la definición de cognoscibilidad como el marco teórico de aprendizaje.

Sea $Z$ un conjunto, $\mathcal{D} = (\Sigma, \mathcal{P})$ distribución sobre $Z$ en lo que sigue.

\begin{definition}[Función de pérdida]
Sea $H$ un conjunto arbitrario, se denomina función de pérdida de $H$ sobre $Z$ a cualquier función de la 
forma:
\[
  l : H \times Z \rightarrow \mathbb{R}_0^{+}
\]
\end{definition}

Modificamos los siguientes puntos de las definiciones básicas en \ref{sec:defs}:

\begin{itemize}
  \item \textbf{Verdadero etiquetado}. No asumimos que exista.
  \item \textbf{Generación de instancias}. Tenemos una distribución $\mathcal{D}$ sobre $Z = (X,Y)$ donde $X$ es el dominio
  y $Y$ el conjunto de etiquetas.
  \item \textbf{Error del clasificador}- Lo redefinimos como:
  \[
    L_{\mathcal{D}}(h) :=  \expect_{z\sim \dist} (l(h,z))
  \]  
  \item \textbf{Error empírico}. Lo redefinimos como:
  \[L_{S} (h) := \frac{1}{m} \sum_{i=1}^m l(h,z_i)\]
\end{itemize}

\begin{definition}[APAC cognoscible]
Una clase $H$ es agnósticamente PAC (APAC) cognoscible respecto a $Z$ y respecto a una función de pérdida 
$l: H \times Z \rightarrow \mathbb{R}^{+}$ si existe una función 
$m_{H} : ]0,1[^2\rightarrow \mathbb{N}$ y un algoritmo 
$A: \underset{m\in \mathbb{N}}{\bigcup} Z^m \rightarrow Y^X$ verificando que si 
$0 < \varepsilon, \delta < 1$, entonces para toda distribución $\mathcal{D}$ sobre $Z$ que haga medible a $z \mapsto l(h, z)$,
se cumple que dado $m\ge m_{H}(\varepsilon, \delta)$ entonces:
\[\mprob \bigg[L_{\mathcal{D}}(A(S)) \le \underset{h\in H}{inf} L_{\mathcal{D}}(h) + \varepsilon \bigg] \ge 1-\delta\]
\end{definition}

Nótese que en contraste con nuestra definición de PAC cognoscibilidad, el algoritmo $A$ podía devolver una 
hipótesis que no estuviera en $H$. Cuando permitimos que el algoritmo $A$ devuelva una función 
$h \notin H$, de manera que $h \in H'$ y  $H \subset H'$ una clase de funciones donde la función de pérdida 
es extensible de manera natural, el aprendizaje recibe el nombre de \textit{aprendizaje impropio}. Si suponemos
$A(S) \in H$ estaremos usando \textit{aprendizaje propio}.

\subsection{Particularización a distintos paradigmas}

\subsubsection{Clasificación binaria}
\label{sec:clas-binaria}

Tendremos $Y=\{0,1\}$ o $Y=\{-1,1\}$ conjunto de etiquetas, un conjunto $X$ y $Z=X\times Y$, y nuestras clases
de hipótesis $H$ estarán contenidas en $Y^X$, de manera que al aprender distribuciones sobre $Z$ podríamos 
tener etiquetas distintas para una misma instancia $x \in X$.

\begin{definition}[Función de pérdida $0-1$]
Llamamos función de pérdida $0-1$ a una función de pérdida de $H\subseteq Y^X$ sobre $Z$, $l$, verificando:
\[
  l(h, (x,y)) = \left\{\begin{array}{ll}
                        1 & \textrm{si } h(x)\neq y\\
                        0 & \textrm{en otro caso}
                        \end{array}\right.
\]
para cualquier $h\in H$, cualquier $(x,y) \in Z$
\label{def:zero-one-loss}
\end{definition}

Consideraremos la función de pérdida $0-1$ en clasificación binaria de ahora en adelante, a no ser que se especifique
lo opuesto.

Con estos conceptos revisitados, podríamos asegurar que la hipótesis que menor error comete para 
clasificación binaria y la función de pérdidad $0-1$ es el llamado \textit{clasificador de Bayes}:
\[
  f_{\mathcal{D}}(x) = \left\{\begin{array}{ll}
                              1 & P [y = 1 |x] >= 0.5\\
                              0 & \textrm{caso opuesto}
                              \end{array}\right.
\]

Pero evidentemente asumimos que el algoritmo no tiene acceso a la distribución, sino solo a los datos 
preetiquetados de $S$, y que por tanto no puede devolver el clasificador de Bayes directamente.

Particularizamos a continuación la definición de APAC anterior a clasificación binaria
con función de pérdida $0-1$ y aprendizaje impropio (ya que la función de pérdida $0-1$ está definida 
en $2^X$).

\begin{definition}[APAC cognoscible - paradigma binario]
Una clase de funciones $H \subseteq 2^X$ es APAC cognoscible sii existen 
$m_{H} : ]0,1[^2\rightarrow \mathbb{N}$ y un algoritmo $A: \underset{m\in \mathbb{N}}{\bigcup} Z^m \rightarrow 2^X$ verificando que para todo
$0 < \varepsilon, \delta < 1$ y para toda distribución $\mathcal{D}$ sobre $Z$, dado $m \ge m_H(\varepsilon, \delta)$:
\[
  \mprob \bigg[L_{\dist,f}(A(S)) \le \inf_{h\in H} L_D(h) + \varepsilon \bigg] \ge 1-\delta 
\]
\end{definition}

\subsection{Clasificación multiclase}
Dado $Y$, verificando $|Y| > 2$, $Z=X\times Y$ y $H\subseteq Y^X$, al problema lo llameremos de 
clasificación multiclase. 

Una posible función de pérdida para este problema sería también la función de pérdida $0-1$.

\subsection{Regresión}
Tenemos un problema de regresión si tomamos $Z=X \times \mathbb{R}$ y $H\subseteq \mathbb{R}^X$
y $l(h,(x,y)) = (h(x)-y)^2$ función de pérdida cuadrática.

\subsection{k-clustering}
Tenemos un problema de clústering si $Z=\mathbb{R}^d$ y los elementos de $H$ son funciones constantes 
$h = (c_1, \ldots c_k)$ con $c_i \in \mathbb{R}^d$. Una posible función de pérdida podría ser 
$l((c_1, \ldots c_k), z) = \min_{i=1,\ldots, k} \norm{c_i - z}^2$.


\section{Relación entre APAC y PAC cognoscibilidad}
Aunque pueda parecer contraintuitivo, las clases APAC cognoscibles están contenidas en las PAC cognoscibles. De hecho, 
tenemos una relación tal cual se muestra en la figura siguiente:

\imgcaption{./imgs/clases-pac.png}{Relación entre clases PAC}{0.6}

\begin{fact}
 Si $H\subseteq 2^X$ es APAC propiamente cognoscible con funciones de pérdida $0-1$, entonces $H$ es PAC cognoscible.
 \label{fact:rel-pac-apac}
\end{fact}

  \begin{proof}
   Sea $f\in H$ y una distribución $\dist=(\Sigma,P)$ sobre $X$ verificando que $H\subseteq \Sigma$. Podemos ver 
   $\dist$ como una distribución que asigna a cada $x\in X$ la etiqueta $f(x)$. Por definición $\inf_{h\in H} L_{\dist}(h) = 0$. 
   Para terminar, dado $S=(z_1, \ldots, z_m)$ con $z_i = (x_i,y_i)$:
   \[
     L_{\dist}(h) = \expect_{x\sim \dist}\bigg(l \big(h,(x,f(x)) \big)\bigg) = \prob[f \cap h] = L_{\dist,f}(h)
   \]
   donde $[f\cap h] \in \Sigma$ por hipótesis ($f,h\in H \subseteq \Sigma$).
   
   Por otro lado:
   \[
     L_{S}(h) = \frac{1}{m} \sum_{i=1}^m l(h,z_i) = \frac{|\{i\in\{1,\ldots, m\}: h(x_i) \neq f(x_i)(= y_i)\}|}{m}
   \]
   
   Además $l(h,\cdot)$ es $\Sigma$ medible ya que $l(h, \cdot)^{-1}(\{0\}) = [f = h] \in \Sigma$ y 
   $l(h, \cdot)^{-1}(\{1\}) = [f\neq h]$. Por hipótesis existiría un algoritmo $A$ y $m_H(\epsilon, \delta)$
   tal que para todo $m\ge m_H(\epsilon, \delta)$:
   \[
     \mprob \bigg[L_{\mathcal{D}}(A(S)) \le \varepsilon \bigg] \ge 1-\delta
   \]
   
   Luego $H$ sería PAC cognoscible.
  \end{proof}
\chapter{Condiciones suficientes para ser PAC cognoscible}
  \section{Clases de Glivenko-Cantelli}

Empezamos definiendo lo que es un conjunto de entrenamiento $\epsilon$ representativo.

\begin{definition*}
 Un conjunto de entrenamiento $S$ es $\epsilon$-representativo de una clase de hipótesis $H$ respecto a una
 distribución $\mathcal{D}$ y una función de pérdida $l$, si $\forall h\in H$ se tiene:
 
 \[|L_S(h) - L_{\dist}(h)| \le \epsilon\]
\end{definition*}

Intuitivamente, cualquier $ERM_H$ sobre $S$ sería un buen algoritmo de aprendizaje 
si $S$ es $\epsilon$ representativo para cualquier $\epsilon \in ]0,1[$ y cualquier distribución. 
Lo formalizamos con la siguiente proposición:

\begin{fact}
 Si $S$ es $\epsilon$-representativo de $H$ con respecto a una distribución $\mathcal{D}$ y una función de 
 pérdida $l$, entonces se tiene:
 
 \[L_{\dist}(ERM(S)) \le \inf_{h\in H} L_{\dist}(h) + 2\epsilon\]
 
 \label{fact:epsilon-rep}
\end{fact}

\begin{proof}
 Sea $A$ un $ERM_H$
 
 $L_{\dist}(A(S)) \le L_S(A(S)) + \epsilon$ por ser $S$ $\epsilon$-representativa. Pero por definición de 
 $ERM_H$, $\inf_{h\in H} L_S(A(S)) + \epsilon \le \inf_{h\in H}L_S(h) + \epsilon$ y aplicando $\epsilon$-representatividad de nuevo:
 
 \[\inf_{h\in H} L_S(h) + \epsilon \le \inf_{h\in H} L_{\mathcal{D}}(h) + 2\epsilon\]
 
 En resumen: $L_{\dist}(A(S)) \le \inf_{h\in H} L_{\mathcal{D}}(h) + 2\epsilon$, para $A \in ERM_H$ arbitrario.
\end{proof}


\begin{definition*} \textbf{Clase de Glivenko-Cantelli}

Decimos que una clase de hipótesis $H$ es de Glivenko-Cantelli, respecto a un dominio $Z$, y a 
una función de pérdida $l$, si existe una función ${m_{H}^{CU}: ]0,1[^2 \rightarrow \mathbb{N}}$ 
verificando que para todo $0 < \delta, \epsilon < 1$ y para toda distribución $\dist$ sobre $Z$, siendo 
$S\sim \dist^m$ con $m \ge m_{H}^{CU}(\epsilon, \delta)$, 
entonces:

\[\mprob [\forall h\in H, |L_S(h) - L_{\dist}(h)| \le \epsilon] \ge 1-\delta\]
\end{definition*}

Análogamente a lo que ocurría en la definición \ref{def:pac}, podemos considerar $m_H^{CU}$ única, en el sentido de
que es la mínima función que satisface las hipótesis.

\section{Glivenko-Cantelli y APAC cognoscibilidad}

\begin{theorem*}
Sea $H$ una clase de hipótesis de Glivenko-Cantelli, respecto a $\mathcal{D}$ y función de pérdida $l$. 
Entonces es APAC cognoscible con cualquier algoritmo $ERM_H$ y complejidad muestral
$m(\epsilon, \delta) \le m_{H}^{UC} \left(\frac{\epsilon}{2}, \delta \right)$ 
\end{theorem*}

  \begin{proof}
  Sea $A$ un $ERM_H$ arbitrario, y $\dist$ una distribución arbitraria sobre $Z$.
  Fijamos $m = m_{H}^{UC} \left(\frac{\epsilon}{2}, \delta \right)$.

  Sea $S = (z_1, \ldots z_m)$ un conjunto de entrenamiento, verificando que: 

  \begin{equation}
    \forall h\in H, |L_{S}(h)-L_{\dist}(h)| \le \frac{\epsilon}{2}
    \label{eq:gc-ineq}
  \end{equation}

  Entonces $S$ es $\epsilon$ representativa para $\dist$ y $l$, y por proposición \ref{fact:epsilon-rep}:

  \begin{equation}
   L_{\dist}(A(S)) \le \inf_{h\in H} L_{\dist}(h) + \epsilon
   \label{eq:gc-result}
  \end{equation}

  Pero como $\eqref{eq:gc-ineq}$ ocurre con probabilidad (sobre $S$) mayor o igual a $1-\delta$, entonces $\eqref{eq:gc-result}$
  ocurre con probabiliad mayor o igual a $1-\delta$
  
  \end{proof}
  
  
\begin{fact}
Sea $H$ una clase de hipótesis finita, $Z$ un dominio y sea $l : H \times Z \rightarrow [a,b]$ una función de pérdida. Entonces $H$ verifica la propiedad de convegencia uniforme con: 

\[m_{H}^{CU}(\epsilon, \delta) \le \left\lfloor \frac{log(2|H|/\delta)(b-a)^2}{2\epsilon^2} \right\rfloor + 1\]
\end{fact}
  \begin{proof}
  Sea $H$ una clase de hipótesis finita.

  Fijamos $0 < \delta, \epsilon < 1$. Necesitamos encontrar $m\in \mathbb{N}$ verificando:

  \[P_{S\sim \mathcal{D}^m} [\exists h\in H |L_S(h) - L_{\mathcal{D}}(h)| > \epsilon] < \delta\]

  Partimos de la siguiente desigualdad, que usaremos más adelante, obtenida por subaditividad:

  \[P [\exists h\in H |L_S(h) - L_{\mathcal{D}}(h)| > \epsilon] \le \sum_{h \in H} P [|L_S(h) - L_{\mathcal{D}}(h)| > \epsilon]\]
  Fijamos $h \in H$.

  Dado un conjunto de entrenamiento $S_z = (z_1, \ldots z_m)$, recordamos que $L_{\mathcal{D}} (h) = \mathbb{E}_{z\sim \mathcal{D}} [l(h,z)]$ y que $L_{S_z}(h) = \frac{1}{m} \sum_{i=1}^m l(h,z_i)$

  Donde $z_i \sim \mathcal{D}$ y por tanto $\mathbb{E}_{S \sim \mathcal{D}^m} [L_S(h)] = \mathbb{E}_{z \sim \mathcal{D}} [l(h,z)] = L_{\mathcal{D}} (h)$. Además, llamando $X_i = l(h,Z_i)$, por ser $S=(Z_1, \ldots Z_m)$ m.a.s que genera los conjuntos de entrenamiento, se tiene que las $X_i$ son independientes e idénticamente distribuidas, con $P[a \le X_i \le b] = 1$. Estamos en condiciones de aplicar la desigualdad de Hoeffding.

  Por tanto:

  \[P \left[\left| \frac{1}{m} \sum_{i=1}^m X_i - L_{\mathcal{D}} (h) \right| > \epsilon\right] = P [|L_S(h) - L_{\mathcal{D}}(h)| > \epsilon] \le 2e^{-2m \left( \frac{\epsilon}{b-a} \right)^2}\]

  Y por tanto:

  \[P [\exists h\in H |L_S(h) - L_{\mathcal{D}}(h)| > \epsilon] \le |H| 2e^{-2m \left( \frac{\epsilon}{b-a} \right)^2}\]

  Despejando $m$ para que $|H| 2e^{-2m \left( \frac{\epsilon}{b-a} \right)^2} < \delta$ llegamos al resultado buscado.
  \end{proof}

Recordemos hasta ahora el resultado que habíamos obtenido era su carácter PAC cognoscible, donde 
agnósticamente PAC cognoscible y cognoscible con funciones de pérdida 0-1 era un término equivalente. El 
teorema que enunciamos a continuación, deducible a partir del teorema sobre el caracter agnóstico - PAC 
cognoscible de clases de funciones con propiedad de convergencia uniforme, en particular las finitas, 
generaliza el resultado para cualquier funciones de pérdida acotada.

\begin{theorem*}

Sea $H$ una clase de hipótesis finita, $Z$ un dominio y sea $l : H \times Z \rightarrow [a,b]$ una función de pérdida. Entonces $H$ es PAC cognoscible con complejidad muestral:

\[m_{H}( \epsilon, \delta ) \le \left\lceil \frac{2 log(2|H|/\delta)(b-a)^2}{\epsilon^2} \right\rceil\]
\end{theorem*}
  \begin{proof}
  Es trivial desde el anterior teorema y el hecho de que convergencia uniforme implica ser agnósticamente PAC cognoscible
  \end{proof}

  
% *************************************************************************************************************
% Parte de informática
% *************************************************************************************************************
\part{Informática}

\bibliography{references}
\bibliographystyle{IEEEtran}

\end{document}